\chapter{A Network That Learns by Rewriting}
\label{ch:spiking}

The four preceding chapters built learners at four sizes: the individual with its
budget, the game it plays, the composite that a population of individuals is, and the
enzyme that individuates one from another. Each was argued at the scale where the
argument is interesting and none was small enough to watch run.

This chapter is the small one. It takes the weighting construction of
Chapter~\ref{ch:weight} and points it at the smallest system in which the whole
apparatus is visible at once --- rate constants, multiplicity, plasticity, and the
funding gate --- and what comes out is a spiking neural network with learning in it. The
network is not an analogy and it is not an encoding. It is what
Definition~\ref{wt:def:propensity} computes when the theory is a handful of
for-comprehensions.

The reason it belongs in this part rather than beside the construction is the last of
those four factors. A neuron here has an energy budget, because the funding gate of
Proposition~\ref{wt:prop:gate} applies to it like anything else, and a neuron whose local
supply is exhausted is silent regardless of how strong its synapses are. Weight is what a
synapse is worth; cost is what a spike costs. That a strong synapse on a starving neuron
transmits nothing is a statement neither decoration makes alone, and it is the mortal
computation of \chSci\ arriving at a scale where one can count the tokens.

\section{The design decision that makes the example an example}

One can always model a neural network in a process calculus by writing the arithmetic
into the payloads: send a real number down a channel, multiply it by a stored weight,
sum, squash, forward. Nothing in Chapter~\ref{ch:weight} would be doing any work; the
calculus would be a programming language and the network a program written in it.

We do the opposite. The synaptic efficacies are \emph{not} in the terms. They are the
values of the weight map, keyed by formulae picking out the synapses. In one sentence:

\begin{quote}
\emph{Synaptic efficacy is represented as transition intensity in the operational
semantics, and the same transition updates that intensity.}
\end{quote}

Three consequences follow, and the third is the one the chapter exists for. Firing is a
stochastic event rather than an arithmetic one. The network's response nonlinearity is a
property of the chain rather than a function written down anywhere. And synaptic
plasticity is literally the update functions of Definition~\ref{wt:def:augmented}, so
that a network which learns and a network which infers are the same weighted theory
running the same transition relation.

It is worth being careful about what the object is, since the temptation to overname it
is strong. It is an asynchronous stochastic spiking recurrent network with plastic
transition rates. It is not a recurrent neural network in the machine-learning sense:
there is no backpropagation, no membrane potential, no refractory period, no global clock
and no discrete timestep, and the quantity compared below with a textbook neuron's
pre-activation is a hazard rate rather than an affine form awaiting an activation
function. Nothing in the argument needs the stronger name, and the claim it does support
is the more interesting one.

\section{Neurons are persistent for-comprehensions}

\begin{definition}[Neuron]
\label{spk:def:neuron}
Let neuron $i$ have presynaptic channels $a_{1i}, \ldots, a_{k_i i}$ --- its dendrites ---
and postsynaptic targets $\mathrm{post}(i)$. Its \emph{threshold presentation at
$\theta_i$} is the persistent for-comprehension
\[
  N_i \;\triangleq\;
  \texttt{for}\bigl(\,
    \underbrace{x_{s_1} \Leftarrow a_{s_1 i} \;\&\; \cdots \;\&\; x_{s_\theta} \Leftarrow a_{s_\theta i}}
    _{\text{one group per }\theta_i\text{-subset } S}
  \;;\; \cdots \,\bigr)
  \;\bigl\{\; \textstyle\prod_{j \in \mathrm{post}(i)} a_{ij}!(\mathrm{spk}) \;\bigr\},
\]
with one join group per $\theta_i$-element subset of the dendrites, $\&$ joining within a
group and $;$ separating groups. The network is
$\mathsf{Net} \triangleq \texttt{new}\ \vec{a}\ \texttt{in}\ \{\prod_i N_i \mid M_0\}$,
with $M_0$ an initial multiset of pending spikes.
\end{definition}

Three points, in increasing order of importance.

The receipts are \emph{persistent}. A linear receipt evaporates after one firing, so a
network built from linear receipts models a single sweep and not a network. Persistence
is what makes the object a standing structure rather than a trace through one. The
pending spikes are still consumed per firing; it is the neuron that survives.

Recurrence is not an additional construct. It is the statement that some $a_{ij}$ for
$j \in \mathrm{post}(i)$ feeds a dendrite of a neuron from which $i$ is reachable, $i$
itself included. The unrolled network in the machine-learning sense is the reduction
graph; the hidden state at a given moment is the \emph{marking}, the standing multiset of
pending spikes on the restricted channels. Since the marking is part of the term and the
term is part of the configuration, hidden state, weights and time are carried by one
object. That is Definition~\ref{wt:def:config} paying for itself.

And the threshold is in the join structure.

\subsection{The keys, and why they form a partition}

\begin{definition}[Synaptic keys]
\label{spk:def:keys}
For each synapse $(j,i)$ let $\varphi_{ji}$ be the structural formula holding of a
communication redex on the channel $a_{ji}$ --- an input on $a_{ji}$ separated from an
output on $a_{ji}$ --- and put
$\Phi_{\textsc{comm}} = \{\varphi_{ji}\} \cup \{\varphi_{\bot}\}$ with $\varphi_{\bot}$
the default class of Remark~\ref{wt:rmk:default}.
\end{definition}

\begin{proposition}
\label{spk:prop:admissible}
$\Phi_{\textsc{comm}}$ is admissible in the sense of Definition~\ref{wt:def:reqs}, with
depth bound $\alpha = 2$.
\end{proposition}

\begin{proof}
Distinct synapses are distinct channels and a communication redex is on exactly one
channel, so (P1) holds; $\varphi_{\bot}$ gives (P2) by Remark~\ref{wt:rmk:default}. Each
key has depth $2$, and checking it against a term is a lookup in the bag, so
\ref{wt:req:dec}, \ref{wt:req:loc} and \ref{wt:req:str} hold.
\end{proof}

The value $\dec(\varphi_{ji})$ is the efficacy of synapse $(j,i)$. Nothing else in the
presentation mentions it.

\subsection{Threshold is join structure}

\begin{proposition}[McCulloch--Pitts, as an expressivity lemma]
\label{spk:prop:mcp}
The threshold presentation of Definition~\ref{spk:def:neuron} realizes the monotone
Boolean threshold function ``at least $\theta_i$ of the dendrites carry a pending
spike''. It uses $\binom{k_i}{\theta_i}$ join groups.
\end{proposition}

\begin{proof}
A group is enabled exactly when every channel it binds carries a spike, since $\&$
requires simultaneous availability of all its binds; the receipt is enabled exactly when
some group is, since $;$ separates alternatives. Some $\theta_i$-subset is fully
populated exactly when at least $\theta_i$ dendrites are.
\end{proof}

So the historical origin of the neuron model --- McCulloch and Pitts' threshold
unit~\cite{mcp} --- is not encoded in a weighted theory but \emph{is} a
for-comprehension, with the threshold appearing as the arity of the joins and the
alternatives as the group separator. The qualification ``expressivity lemma'' is not
cosmetic: $\binom{k}{\theta}$ groups is unusable at realistic fan-in, so the scalable
presentation is the one at $\theta = 1$, where threshold behavior is trivial and the
join structure carries no weight. What the lemma establishes is that no apparatus beyond
the for-comprehension is needed to \emph{express} a threshold. What it does not establish
is that this is how one would build a large network.

\begin{remark}[Coincidence versus rate]
\label{spk:rmk:dial}
$\theta_i = k_i$ gives a pure coincidence detector: the neuron fires only on simultaneous
arrival at every dendrite. $\theta_i = 1$ gives a pure rate integrator: any arrival can
trigger it, and the weights determine how fast. Real cortical neurons sit between, and so
does the presentation, with $\theta$ as the dial. This is a rare case in which a
biological parameter and a syntactic parameter are the same parameter.
\end{remark}

\section{The weighted sum, and where the squashing function went}
\label{spk:sec:squash}

Take $\theta_i = 1$, so $N_i$ has $k_i$ single-bind alternatives, and let the marking
assign $m_{ji}$ pending spikes to dendrite $a_{ji}$.

\begin{proposition}[Propensity of a neuron]
\label{spk:prop:sum}
With every geometric factor $1$ and every redex funded, the total propensity of the
redexes belonging to neuron $i$ is
\[
  \prop_i \;=\; \sum_j \dec(\varphi_{ji}) \cdot m_{ji} .
\]
\end{proposition}

\begin{proof}
The receipt is persistent, so exactly one input is available on each $a_{ji}$, and the
marking supplies $m_{ji}$ outputs; the number of redexes classified $\varphi_{ji}$ is
therefore $m_{ji}$. Definition~\ref{wt:def:propensity} gives $\dec(\varphi_{ji})m_{ji}$
per synapse, and the classes are disjoint by Proposition~\ref{spk:prop:admissible}, so
they may be summed.
\end{proof}

Formally this is the weighted sum $\sum_j w_{ji}m_{ji}$ of the textbook artificial
neuron, with presynaptic activity read as the marking --- and it has not been written
into any term. It is what Definition~\ref{wt:def:propensity} computes.

The comparison should nevertheless be made carefully, because the two quantities are not
the same kind of thing. $\prop_i$ is a hazard rate, with units of inverse time, governing
a race between exponential clocks. It is not an affine pre-activation waiting to be
passed through an activation function. What the next result shows is that a saturating
nonlinearity nevertheless appears, and appears without being supplied.

\begin{corollary}[The response saturates and nobody wrote it]
\label{spk:cor:saturate}
Hold the marking on $i$'s dendrites fixed over a window $[0,\Delta]$. The probability
that neuron $i$ fires at least once in the window is
\[
  \Pr[\,i \text{ fires}\,] \;=\; 1 - \exp\bigl(-\Delta \textstyle\sum_j \dec(\varphi_{ji})m_{ji}\bigr),
\]
a monotone function of the weighted sum, with slope $\Delta$ at the origin and asymptote
$1$.
\end{corollary}

\begin{proof}
Under a marking held fixed the propensity $\prop_i$ is constant, so firings of $i$ within
the window are the events of a homogeneous Poisson process of rate $\prop_i$, by
Theorem~\ref{wt:thm:markov} restricted to the sub-chain along which the marking on $i$'s
dendrites does not change. The probability of at least one event in $[0,\Delta]$ is
$1 - e^{-\prop_i \Delta}$.
\end{proof}

Two provisos, one standard and one easy to miss. The standard one is that the marking is
not in fact frozen, since other neurons are firing into it; the corollary is a statement
about a window short enough that the input has not moved, which is the usual
justification for a rate-coded neuron and the usual approximation. The one easy to miss
is that $i$'s \emph{own} firings change the marking, since a firing consumes the spike it
fired on. The hypothesis must be read as freezing the marking against all sources of
change, $i$'s own consumption included, which is why the statement is about the first
event rather than about a count.

\begin{remark}[Saturating, not logistic]
\label{spk:rmk:notlogistic}
$1 - e^{-\Delta x}$ is an exponential saturation function and not the logistic
$\sigma(x) = (1+e^{-x})^{-1}$; the distinction is worth keeping because the claim is
worth as much either way and should be stated exactly. \emph{A saturating nonlinear
activation emerges from exponential race semantics without being programmed into the
term.} No squashing function was supplied and none is needed. A neuron under enormous
drive still cannot fire more than once per event, so its firing probability per window is
bounded by $1$ however large the weighted sum grows, and that bound is the waiting-time
law of the chain rather than a modeling choice.
\end{remark}

\subsection{Boundedness}

\begin{lemma}[Boundedness]
\label{spk:lem:bounded}
Let $f_i = |\mathrm{post}(i)|$ be the fan-out of neuron $i$. If $f_i \leq \theta_i$ for
every $i$, the total number of pending spikes is non-increasing along every run and the
reachable term set is finite.
\end{lemma}

\begin{proof}
Firing $N_i$ consumes $\theta_i$ pending spikes, one per bind in the selected group, and
produces $f_i$, so the total marking changes by $f_i - \theta_i \leq 0$. The neurons are
persistent and the channel set is fixed, so a reachable term is determined by its
marking, and markings with at most $|M_0|$ tokens over finitely many channels are finite
in number.
\end{proof}

The hypothesis forbids fan-out exceeding threshold, which most networks want. The repair
we take is a \emph{capacity} guard: bound the occupancy of a target channel after the
firing, which is a saturating synapse and is what a real one does. The marking is then
bounded per channel and the reachable state space is finite unconditionally. It matters
that the bound is read \emph{after} the firing rather than before; the obvious
alternative, refusing a send when the target is already full, blocks the consumer and
deadlocks the network. The other standard repair --- a leak rule discarding a pending
spike at its own rate --- does not bound the state space but makes the marking
positive-recurrent, which is enough for stationary analysis and for classical simulation.

\section{Plasticity is the update function}

\begin{definition}[Local potentiation entry]
\label{spk:def:hebb}
Equip the communication rule's entry at $\varphi_{ji}$ with
\[
  u_{ji}(\dec) \;=\; \dec\bigl[\,\varphi_{ji} \mapsto \min\{\,w_{\max},\ \dec(\varphi_{ji}) + \eta\,\}\,\bigr]
\]
for a potentiation increment $\eta > 0$ and ceiling $w_{\max}$.
\end{definition}

The ceiling is not cosmetic. With a quantized increment it is what makes the reachable
weight-map set finite, which is the hypothesis Definition~\ref{wt:def:hilbert} requires,
and it is a case where a modeling convenience and a mathematical necessity happen to be
the same stipulation.

\begin{theorem}[Learning and inference are one relation]
\label{spk:thm:one}
The transition relation of Construction~\ref{wt:con:functor} on $\Cfg$ has two marginals:
the projection to terms is the network's inference dynamics, and the projection to weight
maps is its learning dynamics. Neither is a separate process. In particular there is no
training schedule, no separate learning clock and no distinguished learning phase, and
the learning rate $\eta$ is not a hyperparameter of an outer loop but a component of the
theory, sitting in the same refinement entry as the weight it modifies.
\end{theorem}

\begin{proof}
Immediate from Definition~\ref{wt:def:config} and Construction~\ref{wt:con:functor}: a
single step produces $(P',\dec')$ with $P'$ obtained by the rewrite and $\dec'$ by the
fold of the update function.
\end{proof}

This is the claim the chapter exists to make, and it is exactly what a static-rate
calculus cannot express. In the Stochastic Pi Machine a rate is fixed at declaration, so
a stochastic $\pi$ model of a neural network can model the firing but not the learning:
one must step outside the calculus, adjust the rates, and re-run. Here the adjustment is
a transition of the same system. Theorem~\ref{wt:thm:spim} says precisely which
restriction is being lifted.

\begin{remark}[How Hebbian this is]
\label{spk:rmk:hebbian}
Definition~\ref{spk:def:hebb} potentiates $\varphi_{ji}$ when a communication on $a_{ji}$
fires --- that is, when a presynaptic spike is consumed by an enabled postsynaptic
receipt. The event is local and activity-dependent, and the condition ``pre and post were
jointly involved'' is not a correlation computed by an observer but the firing of a
single rewrite, which by construction requires both parties. That much is the structural
content of Hebb's proposal~\cite{hebb}, and it is the part worth having: the event which
transmits is the event which potentiates, because a communication is one rewrite.

It is nevertheless less than Hebbian learning as usually meant. Hebbian plasticity
invokes a correlation between pre- and \emph{postsynaptic firing}, and successful
transmission into an enabled receipt is not by itself postsynaptic firing, since the
receiving neuron may not reach threshold. What we have is a minimal local potentiation
rule, and the timing-dependent version costs something specific.
\end{remark}

\begin{remark}[What spike-timing dependence costs]
\label{spk:rmk:stdp}
Spike-timing-dependent plasticity makes the weight change a function of the interval
between pre- and postsynaptic spikes~\cite{stdp}. The ingredients are available, since
the simulator produces a timestamp at every step, but the update function of
Definition~\ref{wt:def:augmented} cannot see them: its type is too narrow. Widening it to
take additionally the matched substitution, the position, the current time and a bounded
\emph{trace summary} --- for this purpose, one decaying eligibility trace per synapse ---
suffices. Theorem~\ref{wt:thm:markov} survives the widening provided the trace summary is
finite-dimensional and updated deterministically, since one simply enlarges the
configuration to carry it. It does \emph{not} survive an update function reading unbounded
history, which would make the process genuinely non-Markovian and put exact simulation out
of reach. The eligibility-trace formulation is precisely the standard device for staying
inside the bounded case, and it is pleasant that the constraint arrives here from the
requirement that the chain be a chain rather than from computational convenience.
\end{remark}

\section{Uniformity: one theory for every network}
\label{spk:sec:uniform}

Definition~\ref{spk:def:keys} gives one key per synapse, so a network with $10^4$
synapses has $10^4$ keys and a network that grows needs a new theory. Reflection removes
both problems. Take the synapse channels to be \emph{structured names} built from the
network's own name and the two endpoint indices, so that a name predicate can recover the
endpoints. Then
\[
  \varphi_{\mathrm{syn}} \;\triangleq\;
  \exists j, i.\ \text{a communication redex on the synapse channel of } (j,i)
\]
is a \emph{single} key describing the whole \emph{namespace} of this network's synapses,
which is namespace logic doing what it was built for~\cite{NamespaceLogic}. The weight map
then has one entry per network rather than per synapse, and per-synapse efficacies are
recovered by letting the entry's value be a function of the recovered indices --- the
standard move of trading a large map for a small map with structured keys.

The point generalizes well beyond bookkeeping, and it is the reason this chapter sits in
this part of the book. A namespace is the computationally coherent notion of scope
(Chapter~\ref{ch:scope}), and what the reflective calculus buys here is that a
\emph{population} of synapses can be priced by one formula because the population is a
nameable thing. In an atomic-name calculus none of this is available: the indices would
have to travel as message payloads, and the logic could reach them only by observing
traffic. This is the same manoeuvre by which \chComp\ takes a learner to be a weighted
population and hence a namespace, run at the scale of a synapse instead of the scale of a
scientific community.

\section{The complex instance, and an honest negative}

Replace the efficacies by amplitudes, keeping the keys, the neurons and the capacity
bound. Definition~\ref{wt:def:jump} gives a jump operator per synapse and
Definition~\ref{wt:def:hamiltonian} a Lindbladian. The example satisfies the finiteness
hypothesis in both components, which is part of why it is the example: the capacity bound
makes the reachable term set finite and the ceiling $w_{\max}$ with a quantized increment
makes the reachable map set finite. So the space is finite-dimensional over
\emph{configurations}, and the projectors supplied by the generated logic have a reading
--- the projector at $\varphi_{\mathrm{syn}}$ has expectation the probability that the
network has an available synaptic transmission at time $t$.

\begin{remark}[Indistinguishable spikes do not interfere]
\label{spk:rmk:bosonic}
$m$ indistinguishable pending spikes on $a_{ji}$ give $m$ redexes with a common
contractum. It is tempting to conclude that the quantum network is superlinear in
presynaptic multiplicity, transmitting at $m^2$ times the single-spike rate where the
classical network transmits at $m$ times. Under the normalization of
Definition~\ref{wt:def:jump} it is not: the amplitude of the synaptic jump channel is
$\sqrt{m}$ times the single-spike amplitude, its weight is $m$ times the single-spike
rate, and at $H = 0$ the quantum network's population dynamics is exactly the classical
network's.

What the complex reading adds is therefore not an enhancement of transmission but a
coherent channel between configurations that the equational presentation declines to
quotient --- and Definition~\ref{spk:def:neuron} quotients everything, so as written it
adds nothing at all. This is worth stating plainly because it is the negative instance of
Principle~\ref{wt:prin:slogan}, and a principle whose negative instances are never
displayed is not doing any work. It also names what would have to be done: construct a
spiking network whose equational theory leaves something coherent, so that
Definition~\ref{wt:def:hamiltonian} has material to work with. We have not done it, and we
do not claim that concurrency supplies it for free.
\end{remark}

\section{What this example does not show}

The learning here is local. Nothing above implements backpropagation, which needs a
global error signal propagated against the direction of inference. One can encode such a
signal as traffic on a second family of channels with its own keys, and the apparatus does
not forbid it, but the result would no longer have the property that makes
Theorem~\ref{spk:thm:one} interesting --- the error channel is an outer loop wearing a
costume. Whether a genuinely local learning rule of this kind can match backpropagation is
not a question answered here, and the honest statement is that the example demonstrates a
\emph{mechanism} for plasticity inside the calculus, not a competitive learning algorithm.

Nor is biological fidelity claimed. The neuron of Definition~\ref{spk:def:neuron} has no
membrane potential, no refractory period and no dendritic geometry. The first two are easy
additions --- a potential is a payload, a refractory period is a guard --- and the third is
what the geometric factor of Remark~\ref{wt:rmk:geometric} is for, which is a use we have
not developed.

What the example does show is worth restating against the four chapters that precede it.
Those chapters argued that a learner is a population with a budget, that its inquiry
terminates because assays are paid for, and that a community of such learners composes.
This one exhibits the same architecture at a scale where every quantity is countable: the
population is a marking, the budget is a funding gate, the inquiry is a race between
exponential clocks, and the learning is the same relation as the inference. It is the
smallest thing in this book that is recognisably a mind-shaped object, and nothing in it
was put there by hand except the rules.

% ===================================================================
